\documentclass[11pt,a4paper]{article}
\usepackage[margin=2.3cm]{geometry}
\usepackage{amsmath,amssymb,booktabs,graphicx}
\usepackage[colorlinks=true,linkcolor=blue,citecolor=blue,urlcolor=blue]{hyperref}
\newcommand{\Rsun}{R_\odot}
\newcommand{\kms}{\,\mathrm{km\,s^{-1}}}
\title{SolarRuns: a flare-triggered Moreton wave with the SurfaceCode MHD solver\\
\large Adopted physical parameters, driver calibration, and references}
\author{SurfaceCode --- SolarRuns}
\date{16 July 2026}

\begin{document}
\maketitle

\section{Purpose and scenario}

This document records the physical parameter values (with sources) adopted for the
\texttt{SurfaceCode/SolarRuns} simulation of a \emph{Moreton wave}: the arc-shaped
H$\alpha$ front that races away from large flares at $500$--$2000\kms$
\cite{Moreton1960,MoretonRamsey1960,Warmuth2004a}. Following Uchida \cite{Uchida1968}, the
front is modelled as the chromospheric footprint of a \emph{coronal fast-mode shock}: a
flare-driven blast propagating through the quiet corona whose lower flank refracts
downward onto the chromosphere (the fast speed increases with height through the low
corona, bending wavefronts back toward the surface). The simulation solves compressible
ideal MHD in a 2D $(x,z)$ vertical slice spanning photosphere $\to$ chromosphere $\to$
transition region $\to$ corona, with the flare modelled as an impulsive-phase pressure
pulse in the low corona --- the ``blast-wave'' scenario
\cite{VrsnakCliver2008,Warmuth2015}, implemented essentially as in the closest published
analogue, the FLASH Moreton-wave simulations of Krause et al.\ \cite{Krause2015}. The
solver is the validated \texttt{MHD2D} core of the SurfaceCode pilot (HLL finite volume,
MUSCL/minmod, SSP-RK2, Dedner GLM divergence cleaning, well-balanced gravity source;
stratified wave propagation verified against the exact linear solution of
\cite{PopescuBraileanu2019} to $<2\%$).

\section{Model atmosphere --- two modes}

\textbf{Production mode (\texttt{MODE=krause}, default).} The background of the published
Moreton-wave simulations of Krause et al.\ \cite{Krause2015}: a \emph{gravity-free,
isobaric} corona ($T=1.6$ MK, $p=2.65\times10^{-3}$ Pa, $n\simeq1.2\times10^{8}\;
\mathrm{cm^{-3}}$; gravity is negligible below $\sim$100 Mm since the coronal pressure
scale height is $\simeq100$ Mm at 1.6 MK) with a dense chromospheric slab ($T=10^{4}$ K,
density contrast 160) below $z=5$ Mm, joined by a resolvable smoothstep edge. This
background is an \emph{exact static solution at any resolution}, which is what makes the
strong-shock/chromosphere interaction numerically robust --- the reason Krause et al.\
adopted it.

\textbf{High-resolution mode (\texttt{MODE=valc}).} The full hydrostatic stratification,
$dp/dz=-\rho g_\odot$ with constant $g_\odot=274\;\mathrm{m\,s^{-2}}$, a piecewise
VAL-C temperature profile \cite{VAL1981}, and a smooth mean molecular weight $\mu(T)$
interpolating between neutral ($\mu=1.25$) and ionized H+He ($\mu=0.6$). Physically
complete (photosphere, temperature minimum, chromosphere, TR, corona --- the full Uchida
waveguide), but hydrostatics fixes the chromosphere$\to$corona density contrast at
$\sim$$10^{5}$, so the density scale length at the transition region is
$\sim$0.2--0.3 Mm for \emph{any} resolvable $T(z)$ path; surviving the blast impact
requires $\Delta z\lesssim75$ km (a multi-hour run; the $\Delta z\ge100$ km attempts
consistently failed with TR-cell positivity losses at the flare column). Adopted anchor
values (VAL III, Table 12, model C, unless noted):

\begin{table}[h]\centering\small
\begin{tabular}{llll}
\toprule
quantity & adopted value & VAL-C / literature value & ref.\\
\midrule
$T$ at $\tau_{500}{=}1$ ($z{=}0$) & $6420$ K & $6420$ K (kinetic, \emph{not} $T_{\rm eff}{=}5772$ K) & \cite{VAL1981,Prsa2016}\\
gas pressure at $z{=}0$ & $1.172\times10^{4}$ Pa & $1.172\times10^{5}\;\mathrm{dyn\,cm^{-2}}$ & \cite{VAL1981}\\
temperature minimum & $4170$ K at $z{=}515$ km & $4170$ K at $h{=}515$ km & \cite{VAL1981}\\
chromospheric plateau & rise to $7000$ K at $2.0$ Mm & $6040$--$7160$ K over $1.1$--$2.0$ Mm & \cite{VAL1981}\\
transition region & centred $2.40$ Mm, width $0.25$ Mm & $2.5\times10^4\to4.5\times10^5$ K over & \cite{VAL1981}\\
 & (numerically widened) & $\sim$280 km at $2.27$--$2.54$ Mm & \\
coronal temperature & $1.5$ MK, isothermal & $1.4$--$1.6$ MK quiet Sun & \cite{WarmuthMann2005,Krause2015}\\
coronal density ($z{=}10$ Mm) & $n_{\rm H}\simeq1.9\times10^{8}\;\mathrm{cm^{-3}}$ & $n_e\simeq1.2\times10^{8}$ \cite{Krause2015}; range & \cite{Krause2015,Warmuth2015}\\
 & (hydrostatic result, not tuned) & $1$--$10\times10^{8}\;\mathrm{cm^{-3}}$ & \\
coronal sound speed & $c_s\simeq185\kms$ & $180\kms$ at $1.4$ MK & \cite{WarmuthMann2005}\\
scale heights & photosph.\ $\simeq150$ km; corona & $H_\rho\simeq130$ km; $\simeq47\,{\rm Mm}\times T/{\rm MK}$ & \cite{VAL1981,Aschwanden2005}\\
 & $\simeq70$ Mm (emergent) & & \\
\bottomrule
\end{tabular}
\caption{Atmosphere. The hydrostatic integration through the adopted $T(z)$ reproduces
the canonical quiet-Sun coronal base density without tuning.}
\end{table}

\section{Magnetic field and the fast-speed waveguide}

Uniform, vertical $B_0=3$ G (default). Quiet-Sun coronal field determinations give
$\sim$2--6 G with a best value near 3 G (EIT-wave seismology, \cite{West2011}); Warmuth
\& Mann's global model uses an effective coronal base field of $3.4$ G
\cite{WarmuthMann2005}; the Krause et al.\ Moreton simulations bracket $B_0=1$--$10$ G
uniform \cite{Krause2015}. With $B_0=3$ G and the model corona,
$V_A(10\,{\rm Mm})\simeq430\kms$ --- matching the derived literature check ($B=3$ G,
$n_e=1.2\times10^8\;\mathrm{cm^{-3}}\Rightarrow V_A\simeq430\kms$) --- and the ambient
perpendicular fast speed is $V_f=\sqrt{c_s^2+V_A^2}\simeq470\kms$, consistent with the
$\sim$300$\kms$ base magnetosonic speed rising with height of \cite{WarmuthMann2005} and
with observed EIT-wave speeds ($170$--$465\kms$, mean $271\kms$ \cite{Klassen2000})
sitting at or below it. Lateral propagation across the vertical field is a
\emph{perpendicular} fast wave; the $V_f(z)$ profile has the deep minimum
($\sim$7--10$\kms$) around the temperature-minimum layer and the monotonic rise into the
corona that the Uchida refraction mechanism requires \cite{Uchida1968}. A Moreton front
at $500$--$2000\kms$ is thus Mach $\sim$5--20 relative to the \emph{chromospheric} fast
speed --- the standard argument that it must be coronally driven \cite{Warmuth2015} ---
and Mach $1.3$--$3$ relative to the corona.

\section{Flare driver (``initial pulse'')}

Impulsive-phase heating of a Gaussian low-corona volume,
\begin{equation}
\dot q(x,z,t)=q_0\,
\exp\!\Big[-\tfrac{(x-x_{\rm fl})^2}{2\sigma_x^2}-\tfrac{(z-z_{\rm fl})^2}{2\sigma_z^2}\Big]\,
\sin\!\big(\pi t/\tau_{\rm fl}\big),\qquad 0\le t\le\tau_{\rm fl},
\end{equation}
with $q_0$ set so the time-integrated central heating equals
$A\,p_0(z_{\rm fl})/(\gamma-1)$: a peak overpressure ratio $\Delta p/p_0\simeq A$ before
expansion. The kernel is elliptical because hydrostatic \emph{pressure} at the transition
region rivals the low-corona value: a round $\sigma=5$ Mm kernel at $z_{\rm fl}=10$ Mm
deposits $\Delta p/p\sim6$ \emph{in situ} at the TR and detonates the thin layer from
inside, whereas the wave physics requires the chromosphere to feel only the arriving
coronal shock (Krause et al.\ avoid this by placing the pulse at 35 Mm; the elliptical
kernel achieves the same separation in our 24 Mm box, TR deposit $\Delta p/p\lesssim0.2$).

\begin{table}[h]\centering\small
\begin{tabular}{llll}
\toprule
driver parameter & adopted & literature basis & ref.\\
\midrule
overpressure ratio $A$ & $100$ (krause) / $30$ (valc) & Vr\v{s}nak--Cliver fiducial $20$ ($M_A{\simeq}1.4$), & \cite{VrsnakCliver2008,Krause2015}\\
 & & $40$ for $M_A{\simeq}2$; Krause et al.\ up to $100$ & \\
source size $\sigma_x\times\sigma_z$ & $5\times2.5$ Mm (elliptical) & flare-kernel $x_0\lesssim10$ Mm; pulse region & \cite{VrsnakCliver2008,Krause2015}\\
 & & $10$ Mm in Krause et al. & \\
source height $z_{\rm fl}$ & $10$ Mm & low-corona pulse, well above the TR & \cite{VrsnakCliver2008,Krause2015}\\
 & & ($35$ Mm in \cite{Krause2015}) & \\
duration $\tau_{\rm fl}$ & $60$ s & $40$--$50$ s pulses in \cite{Krause2015} ($40$ s = min.\ for & \cite{Krause2015,VrsnakCliver2008}\\
 & & observable Moreton signature); impulsive & \\
 & & phase $\sim$1--5 min & \\
temperature cap (implied) & peak $T\sim3\times10^{7}$ K & flare-loop limits $T\le4\times10^{7}$ K, & \cite{Krause2015,Aschwanden2005}\\
 & & $n\le10^{11}\;\mathrm{cm^{-3}}$ & \\
wave energy content & $\sim$few$\,\times10^{20}$ J & observed wave energies $10^{19}$--$10^{21}$ J, & \cite{Ballai2005,Warmuth2015}\\
 & (2D slab $\times$ 30 Mm ribbon) & i.e.\ $\sim$0.01--1\% of the $10^{22}$--$10^{25}$ J & \cite{Emslie2012}\\
 & & flare budget & \\
\bottomrule
\end{tabular}
\caption{Flare-pulse calibration. The blast-wave route requires low-$\beta$ AR-periphery
conditions ($\beta_0\sim0.01$--$0.1$) and impulsive lateral expansion
\cite{VrsnakCliver2008,Pomoell2008}.}
\end{table}

\section{Observational targets}

The run is judged against the observed phenomenology:
\begin{enumerate}\itemsep2pt
\item \textbf{Speeds:} first measured front velocity $845\pm162\kms$, event-mean
$643\pm179\kms$ (12-event statistics \cite{Warmuth2004a}); historical range
$500$--$2000\kms$ \cite{Moreton1960,SmithHarvey1971}.
\item \textbf{Deceleration:} mean $\bar a\simeq-1.5\;\mathrm{km\,s^{-2}}$ (weakening with
time), decaying amplitude, broadening profile --- a freely decaying nonlinear fast-mode
shock relaxing toward the ambient speed \cite{Warmuth2004a,Warmuth2004b}.
\item \textbf{Visibility range:} first seen at $\sim$$97\pm26$ Mm from the source, lost at
$\sim$$301\pm47$ Mm \cite{Warmuth2004a}.
\item \textbf{Chromospheric swing:} downward push of $\sim$1--4$\kms$ (max $4.1\kms$
upper chromosphere \cite{Cabezas2019}; $2.6\kms$ Doppler \cite{Balasubramaniam2010}),
down--up relaxation over $\sim$1--3 min.
\item \textbf{Coronal counterpart:} EUV front, compression ratio $1.05$--$1.3$
\cite{Muhr2011}, SXR fast-mode Mach $1.1$--$1.3$ \cite{Narukage2002,Hudson2003}; slower
EIT fronts at $\sim$$1/3$ of the coronal shock speed in the Chen et al.\ picture
\cite{Chen2002,Chen2005}.
\end{enumerate}

\section{Numerical setup and interventions}

2D $(x,z)$ slice. Krause mode: $L_x=250$ Mm, $z\in[0,60]$ Mm, $\Delta x=250$ km,
$\Delta z=250$ km, pulse at $(x,z)=(40,35)$ Mm, $A=100$ (the observable-Moreton amplitude of
\cite{Krause2015}), $B_0=3$ G, $\tau_{\rm fl}=50$ s; H$\alpha$ proxy trace at the upper
slab $z=4$ Mm (where the skirt swing is largest), coronal trace at $z=15$ Mm. VAL-C mode: $L_x=200$ Mm, $z\in[-0.6,24]$ Mm, $\Delta x\le250$ km,
$\Delta z\lesssim75$ km; flare centred at $x_{\rm fl}=30$ Mm;
mass-conserving absorbing layers (velocity $+$ internal-energy relaxation toward the
hydrostatic background) at the top and both lateral boundaries. Diagnostics: 2D snapshots
of $v_z$, $v_x$, $\Delta p/p_0$ every 30 s; distance--time stacks (1 s cadence) of
chromospheric $v_z$ at $z=1.5$ Mm (H$\alpha$ proxy) and coronal $\Delta p/p_0$ at
$z=8$ Mm (EUV proxy); front tracking and kinematics.

\medskip\noindent
\textbf{Known simplifications.}
(i) Ideal MHD: no radiative losses, thermal conduction, or partial-ionization physics ---
acceptable on the minutes-long wave-crossing time, but it leaves the \emph{flare site
itself} with no physical relaxation channel. We therefore gently relax the velocity in a
Gaussian column around the flare site after the pulse ends ($t>\tau_{\rm fl}$, rate
$0.1\;\mathrm{s^{-1}}$, $\sigma_x=2\sigma_{\rm fl}$, $z\lesssim10$ Mm): in ideal MHD the
post-flare cavity refill otherwise steepens indefinitely against the numerically thin
transition region. The escaping front is $\gtrsim V_f\tau_{\rm fl}\sim30$ Mm away when
the relaxer activates; front kinematics are read at $x-x_{\rm fl}\gtrsim30$ Mm.
(ii) The flare heating is applied along the whole column (its low-atmosphere share is
small in relative terms, $\Delta p/p\lesssim0.3$ below the TR, and mimics ribbon
co-heating); test runs that artificially confined the heating to the corona created an
unphysical razor pressure interface on the transition region.
(iii) Uniform vertical field: no arcade geometry, no field-line-stretching (pseudo-wave)
component --- this run isolates the \emph{wave} part of the EIT-wave duality
\cite{Chen2002,Warmuth2015}.
(iv) 2D slab: no cylindrical geometric spreading, so amplitude decay along $x$ is a
conservative (upper) limit on front survival.
(v) A positivity backstop: the HLL scheme is not positivity-preserving, and
blast-driven rarefactions can empty a (numerically thin) transition-region cell; cells
falling below $10^{-3}\rho_0(z)$ are floored with their momentum zeroed (physical wake
evacuation stays $\gtrsim10^{-1}\rho_0$, so the backstop only catches numerical
failures; hit counts are logged).
(vi) The H$\alpha$ signature is represented by the chromospheric velocity swing, not
synthesized radiative transfer (an 8\% optically-thin compression threshold is the
visibility criterion used by \cite{Krause2015}).



\section{Results: comparison with the observed 2006 December 6 Moreton wave}

The production run (\texttt{MODE=krause}, $A=100$, $B_0=3$ G, $\Delta x=\Delta z=250$ km)
is compared against the best-observed recent Moreton wave, the 2006 December 6 event
(3B/X6.5 flare in AR 10930), analysed by Balasubramaniam et al.\ \cite{Bala2010}
(\texttt{RefMoreton/}) --- the same event modelled by Krause et al.\ \cite{Krause2015}.
Front tracking mirrors their method: the leading edge is followed and fitted with the
constant-deceleration model $d=d_0+V_0t+\tfrac12at^2$ (their Eq.~1). Figure
\texttt{plots\_production/fig5\_obs\_comparison.png}:

\begin{table}[h]\centering
\begin{tabular}{lll}
\toprule
 & simulation & observed \cite{Bala2010}\\
\midrule
mean front speed & $857\kms$ (41--226 s) & $\sim850\kms$ (18:43--18:50 UT)\\
constant-deceleration fit $V_0$ & $891\kms$ & $\sim1125\kms$\\
deceleration $a$ & $-0.42\;\mathrm{km\,s^{-2}}$ & $-0.87\;\mathrm{km\,s^{-2}}$\\
final tracked speed & $813\kms$ & $\sim$700--800$\kms$ (approaching ambient)\\
chromospheric swing amplitude & $2.0\kms$ (far field) & 1--4$\kms$ (typical events
\cite{Cabezas2019,Balasubramaniam2010})\\
chromospheric signal range & 31--126 Mm from source & $\sim$120--550 Mm from radiant point\\
coronal front tracked to & 204 Mm ($\Delta p/p<3\%$) & EUV counterpart to $\gtrsim R_\odot$ scales\\
\bottomrule
\end{tabular}
\caption{Simulated vs.\ observed kinematics of the 2006 December 6 Moreton wave.}
\end{table}

\noindent\textbf{Agreement.} The propagation-phase kinematics — the observable Moreton
signature — are reproduced: the mean front speed matches the observed $850\kms$ to within
the measurement error, the front decelerates at the same order ($-0.4$ vs.\
$-0.9\;\mathrm{km\,s^{-2}}$, both fits weakening with time), the final speed approaches
the ambient fast-mode speed from above (the freely decaying weak-shock behaviour of
\cite{Warmuth2004b,AU2011}), and the chromospheric down--up swing amplitude falls inside
the observed 1--4$\kms$ band.

\medskip\noindent\textbf{Differences, and why.} (i) The observed wave stayed traceable to
$\sim$550 Mm; our chromospheric signal drops below the $0.3\kms$ threshold at
$\sim$126 Mm and the coronal front below $\Delta p/p=3\%$ at $\sim$204 Mm (the domain
ends 210 Mm from the source). The 2006 event was extreme (X6.5, $\sim$270$^\circ$
azimuthal span, the strongest field change of $\sim$80 surveyed flares), our pulse is the
\emph{minimum observable-Moreton} amplitude of \cite{Krause2015}, and a 2D slab
under-represents the sustained lateral driving of a real, volumetric CME expansion ---
all three push the simulated visibility range low. (ii) The observed wave was first
detected only at $\sim$120 Mm because flare emission saturates H$\alpha$ closer in; our
inner cutoff (31 Mm) is the flare-column guard zone --- analogous in effect, not in
cause. (iii) $V_0$ is $\sim$20\% low: the fit begins at our first clean detection
($t=41$ s), missing the strongly super-magnetosonic launch phase the observed
back-extrapolation captures. \emph{Driver caveat:} \cite{Bala2010} conclude that for
this specific event the kinematics favour lateral CME-arcade expansion over a pure flare
pressure pulse (the pulse would need $\sim$100 km\,s$^{-2}$ acceleration and
$-10\;\mathrm{km\,s^{-2}}$ deceleration to fit their back-extrapolated launch). Our
blast-driven run cannot arbitrate the launch mechanism --- but once launched, the
propagating fast-mode front forgets its driver, which is exactly what the agreement in
the propagation phase demonstrates.


\section{Referee test: driver-limited vs.\ dissipation-limited range}

A natural objection: if the impulsive run's front fades by $\sim$200 Mm while the real
Moreton wave was seen to $\sim$550 Mm, perhaps the solver simply \emph{over-dissipates} ---
in which case the T Tauri decay-length result (the parent paper) would be suspect for the
same reason. This is testable directly. If the short range is set by the \emph{driver}
(a weak, impulsive pulse in a finite box) rather than by \emph{dissipation}, then driving
the wave harder and for longer, in a larger domain, at \emph{identical} solver, dissipation
physics and resolution, should extend the range toward the observed value.

We therefore added a sustained CME-piston driver (\texttt{DRIVE=piston}): steady energy
deposition at rate $A\,p_0/[(\gamma-1)\tau_{\rm fl}]$ held over $\tau_{\rm drive}=180$ s
(the CME lateral-acceleration phase, cosine-ramped), with the driven lateral extent growing
as $\sigma_x(t)=\sigma_{\rm fl}+V_{\rm pist}\,t$ ($V_{\rm pist}=800\kms$) so the driving edge
tracks outward with the front --- the continuous lateral push of an expanding CME flank
(Chen et al.\ \cite{Chen2002}; Vršnak \& Cliver \cite{VrsnakCliver2008}), as opposed to the
impulsive blast. Domain $L_x=600$ Mm, $A=150$, run to $t=720$ s.

\begin{table}[h]\centering
\begin{tabular}{llll}
\toprule
run (same solver, physics, $\Delta x=\Delta z=250$ km) & mean speed & front reached & driver\\
\midrule
impulsive blast ($A{=}100$, 210 Mm box) & $857\kms$ & 204 Mm & one-shot\\
sustained CME-piston ($A{=}150$, 600 Mm box) & $833\kms$ & \textbf{553 Mm} & continuous\\
observed 2006 December 6 \cite{Bala2010} & $\sim850\kms$ & $\sim550$ Mm & CME (inferred)\\
\bottomrule
\end{tabular}
\caption{The visibility range is set by the driver, not by numerical dissipation.
Figure \texttt{plots\_extreme/fig6\_driver\_test.png}.}
\end{table}

\noindent The sustained-driver front reaches \textbf{553 Mm} --- essentially the observed
$\sim$550 Mm --- while propagating at the same $\sim$830--850$\kms$ mean speed, and its
distance--time track brackets the observed curve across the whole disk. Nothing about the
dissipation changed between the two runs; only the driver and the box did. \emph{The short
range of the impulsive run is driver-limited, not dissipation-limited}, which is the
premise the objection needs and does not have.

\medskip\noindent\textbf{Honest cost of the extreme driver.} The sustained $A=150$ piston
over-drives the chromospheric response: the far-field vertical swing reaches
$\sim$130$\kms$, far above the observed 1--4$\kms$ (the impulsive run gave a realistic
$2\kms$). So the two runs bracket reality --- the impulsive blast reproduces the
\emph{amplitude} but under-ranges; the extreme piston reproduces the \emph{range} but
over-drives the amplitude --- and a real, moderate CME sits between. Neither the
deceleration nor $V_0$ of the piston run should be read as free-decay values: its fit is
contaminated by the 180 s driven phase (hence the inflated $V_0\sim2500\kms$ and
$a\sim-6\;\mathrm{km\,s^{-2}}$). The decisive, driver-independent quantities are the mean
propagation speed and the range, both of which match.

\medskip\noindent\textbf{Bearing on the T Tauri result.} The parent paper's decay length
does not rest on trusting the scheme's dissipation in the first place: the numerical
dissipation there was \emph{measured} (the dedicated numerical-diffusion test), and at the
converged resolution the numerical decay length is $\sim$100$\times$ longer than the
reported physical $L\simeq0.006\,R_\star$, with $L$ drifting only $+15\%$ over a
$\Delta x=50\!\to\!13$ km refinement ladder (converging, not grid-limited). Moreover the
two decays are different physics --- a Mach-$3$ shock plus \emph{ambipolar} friction in
near-neutral gas (T Tauri) vs.\ a weakly nonlinear ideal-MHD wave (solar corona) --- so
there is no shared dissipation knob whose miscalibration would inflate both. The present
solar test removes even the indirect worry: when the solver is asked to carry a fast-mode
front $\sim$550 Mm at the correct speed, it does.


\subsection{Can a single driver hit both the range and the amplitude? A trade-off scan}

Pushing further: the extreme run reproduces the \emph{range} but over-drives the
chromospheric swing ($\sim$130$\kms$ vs.\ observed 1--4$\kms$). Can an intermediate
sustained driver hit both? We scanned the piston amplitude ($A=20,25,50,150$; the $A=25$
case also tripled the driver width, $\sigma_x=15$ Mm, to test a ``wide-and-gentle'' CME):

\begin{table}[h]\centering
\begin{tabular}{lllll}
\toprule
run & $A$ (width) & mean speed & range ($\Delta p/p>3\%$) & far-field swing\\
\midrule
mid20      & 20            & $726\kms$ & 439 Mm & 27$\kms$\\
widegentle & 25 ($\sigma_x{=}15$) & $802\kms$ & 553 Mm & 78$\kms$\\
mid50      & 50            & $778\kms$ & 548 Mm & 63$\kms$\\
extreme    & 150           & $831\kms$ & 553 Mm & 133$\kms$\\
\midrule
\emph{observed} \cite{Bala2010} & --- & $\sim850\kms$ & $\sim550$ Mm & 1--4$\kms$\\
\bottomrule
\end{tabular}
\caption{Range vs.\ chromospheric-swing trade-off across the sustained-driver scan
(Figure \texttt{plots\_tradeoff/fig7\_tradeoff.png}). Every run that reaches the observed
$\sim$550 Mm range has a swing $15$--$30\times$ above the observed band; the
``sweet-spot'' corner (range $\sim$550 Mm \emph{and} swing 1--4$\kms$) is empty.}
\end{table}

\noindent Two robust facts emerge. (i) The \emph{range} saturates at $\sim$550 Mm for any
$A\gtrsim25$ --- it is set by sustained driving and the detection threshold, not by
amplitude, and it does \emph{not} increase further with a stronger driver. (ii) The
\emph{swing} grows with the total deposited energy; widening the driver at fixed $A$
(\texttt{widegentle}) made it \emph{worse}, not better, because a wider Gaussian deposits
more energy. So a single sustained-driver knob cannot reach the observed corner.

\medskip\noindent\textbf{Why the real Sun hits both, and why this strengthens (not weakens)
the argument.} The missing ingredient is \emph{geometric spreading}, which the 2D slab
omits by construction (\S5, simplification iv). On the real Sun the front expands as an
arc, so its \emph{local} amplitude falls as $\sim r^{-1/2}$ (cylindrical) \emph{in addition}
to dissipation; a front seen at 550 Mm is therefore intrinsically weak there --- large
range \emph{and} small local swing. Our 2D slab has no such spreading, so the only way it
can keep a front detectable to 550 Mm is to keep it strong, forcing a large swing. The
trade-off is thus an artifact of the deliberately 2D geometry, in the \emph{conservative}
direction: with spreading restored, the same front would decay \emph{faster} in amplitude,
not slower. This is the same reason the paper flags the 2D decay lengths as \emph{upper
limits} --- and it is the opposite of the ``code over-dissipates'' worry. The dissipation
is, if anything, under-represented here; the wave still crosses the disk because it is
driven, not because the scheme conserves its energy artificially.

\section{Relation to the T Tauri SurfaceCode pilot}

\texttt{SolarRuns} reuses the SurfaceCode solver core unchanged, but the two applications
sit at opposite ends of the same physical family: both simulate an impulsively driven
fast-magnetosonic front sweeping laterally away from a localized surface source on a
magnetized star --- yet in \emph{different atmospheric channels}, and that difference is
the physics.

\begin{table}[h]\centering\small
\begin{tabular}{lll}
\toprule
 & \textbf{T Tauri pilot (parent)} & \textbf{SolarRuns (this work)}\\
\midrule
question & does the accretion-impact wave & does the flare blast reproduce the\\
 & survive ring$\to$pole? ($T=F_{\rm pole}/F_{\rm ring}$) & observed Moreton phenomenology?\\
star & $0.5\,M_\odot$, $2\,\Rsun$ T Tauri, & the Sun, $g_\odot=274\;\mathrm{m\,s^{-2}}$\\
 & $g\simeq34\;\mathrm{m\,s^{-2}}$ & \\
propagation channel & \emph{dense} photospheric layers & \emph{tenuous} corona; the chromosphere\\
 & ($\tau\gtrsim1$ to $\tau\ll1$), the layer & only \emph{prints} the front (Uchida skirt)\\
 & Cranmer's transport assumes & \\
background & real MLT envelope & Krause slab mode (gravity-free,\\
 & (Saha EOS export), hydrostatic, & isobaric) or VAL-C hydrostatic mode;\\
 & isothermal placeholder above $\tau{=}1$ & \S2\\
magnetic field & ring-latitude dipole, $0.1$--$1$ kG & uniform vertical quiet-Sun $3$ G\\
 & (kilogauss star) & \\
driver & accretion-clump impact: boundary & flare pressure pulse: thermal\\
 & velocity piston, $\sim$Mach 3 at the & overpressure $\Delta p/p_0=50$ in the low\\
 & photosphere, $\tau=150$ s & corona, $\tau_{\rm fl}=50$ s\\
key damping physics & nonlinear shock $+$ \emph{ambipolar} & ideal MHD (fully ionized corona;\\
 & \emph{friction} (neutral photosphere; the & shock dissipation only)\\
 & T-Tauri-specific term) & \\
resolution & $\Delta x=50$ km ($H_p\simeq700$ km) & $\Delta x=\Delta z=250$ km (Mm-scale\\
 & & structures; finer than \cite{Krause2015})\\
new numerics & --- (validated core) & per-stage bounded-state protection,\\
 & & flare-column guard/relaxer zones,\\
 & & gravity-free isobaric mode\\
diagnostic & transmitted flux $\Phi(x)$, decay & front kinematics $d(t)$, $v(t)$;\\
 & length $L$ & chromospheric down--up swing;\\
 & & EUV-proxy front\\
outcome & the dense channel \emph{kills} the wave: & the coronal channel \emph{carries} the\\
 & $L\simeq0.006\,R_\star$ ($\sim$8 Mm), buried & front across $\gtrsim$100--200 Mm at\\
 & in ambient convection & $\sim$500--800$\kms$, decelerating\\
\bottomrule
\end{tabular}
\caption{The same solver, the two opposite transport channels.}
\end{table}

\noindent The two results are mutually consistent rather than contradictory: lateral
fast-wave transport survives only in the tenuous exterior (solar corona), not in the
dense surface layers (T Tauri photosphere) --- and even the dense-layer footprints agree
quantitatively, since the T Tauri decay length ($\sim$8 Mm) and the distance over which
the simulated solar \emph{chromospheric} response stays strong ($\lesssim$100 Mm of the
source, \S5) are both small compared to the stellar circumference. This is precisely the
front-vs-energy dichotomy of the parent paper: a coronal front can cross the disk while
the dense-layer energy density it drags along fades within a fraction of a radius. The
solar run is therefore both a demonstration that the SurfaceCode core handles the
classic Moreton problem, and an independent illustration of why the accretion-powered
surface channel on T Tauri stars --- which has no hot corona to ride --- cannot deliver
launch-capable wave energy to the pole.

\begin{thebibliography}{99}\itemsep1pt\small
\bibitem{Moreton1960} Moreton, G.~E. 1960, AJ, 65, 494 --- discovery of H$\alpha$ flare
waves ($\sim$1000$\kms$).
\bibitem{MoretonRamsey1960} Moreton, G.~E., \& Ramsey, H.~E. 1960, PASP, 72, 357.
\bibitem{Uchida1968} Uchida, Y. 1968, Sol.~Phys., 4, 30 --- Moreton waves as the
chromospheric skirt of a coronal fast-mode wavefront refracting downward.
\bibitem{SmithHarvey1971} Smith, S.~F., \& Harvey, K.~L. 1971, in \emph{Physics of the
Solar Corona}, ed.\ C.~J. Macris (Dordrecht: Reidel), 156 --- classic kinematic
statistics ($658\pm159\kms$).
\bibitem{VAL1981} Vernazza, J.~E., Avrett, E.~H., \& Loeser, R. 1981, ApJS, 45, 635 ---
VAL III semi-empirical quiet-Sun atmosphere; model C, Table 12.
\bibitem{Prsa2016} Pr\v{s}a, A., et al. 2016, AJ, 152, 41 --- IAU nominal solar values
($T_{\rm eff}=5772$ K, $\Rsun=6.957\times10^8$ m).
\bibitem{Warmuth2004a} Warmuth, A., Vr\v{s}nak, B., Magdaleni\'c, J., Hanslmeier, A., \&
Otruba, W. 2004a, A\&A, 418, 1101 --- multiwavelength statistics of 12 Moreton events
(speeds, deceleration, visibility range).
\bibitem{Warmuth2004b} Warmuth, A., et al. 2004b, A\&A, 418, 1117 --- Paper II:
amplitude decay and profile broadening (weak-shock signature).
\bibitem{Warmuth2015} Warmuth, A. 2015, Living Rev.~Sol.~Phys., 12, 3 --- review of
large-amplitude global coronal waves: speeds, energies, wave/pseudo-wave synthesis.
\bibitem{WarmuthMann2005} Warmuth, A., \& Mann, G. 2005, A\&A, 435, 1123 --- global
coronal magnetosonic-speed model ($c_s=180\kms$ at 1.4 MK; base $v_{ms}\simeq300\kms$;
effective quiet-Sun base field 3.4 G).
\bibitem{VrsnakCliver2008} Vr\v{s}nak, B., \& Cliver, E.~W. 2008, Sol.~Phys., 253, 215 ---
origin of coronal shocks: flare pressure-pulse (blast) vs.\ CME-piston; quantitative
pulse recipe (kernel $\lesssim$10 Mm, overpressure 10--40, low-$\beta$ requirement).
\bibitem{Krause2015} Krause, G., C\'ecere, M., Francile, C., Costa, A., Elaskar, S., \&
Schneiter, M. 2015, MNRAS, 453, 2799 (arXiv:1505.02723) --- 2D MHD flare-pressure-pulse
Moreton simulations (FLASH): $\Delta p/p\le100$ over 10 Mm, 40--50 s, $B_0=1$--$10$ G,
$T_{\rm cor}=1.6$ MK, $n=1.2\times10^8\;\mathrm{cm^{-3}}$; chromospheric down--up swing
lags the coronal front by 30--50 Mm.
\bibitem{Francile2013} Francile, C., et al. 2013, A\&A, 552, A3 --- 2006 Dec 06 Moreton
kinematics (companion observations to \cite{Krause2015}).
\bibitem{Chen2002} Chen, P.~F., Wu, S.~T., Shibata, K., \& Fang, C. 2002, ApJ, 572, L99 ---
piston-driven coronal shock (Moreton, $773\kms$) + trailing pseudo-wave EIT front
($\sim$$1/3$ speed).
\bibitem{Chen2005} Chen, P.~F., Fang, C., \& Shibata, K. 2005, ApJ, 622, 1202.
\bibitem{Pomoell2008} Pomoell, J., Vainio, R., \& Kissmann, R. 2008, Sol.~Phys., 253, 249
--- Moreton-capable shocks need impulsive lateral expansion, not high final speed alone.
\bibitem{Thompson1998} Thompson, B.~J., et al. 1998, Geophys.~Res.~Lett., 25, 2465 ---
discovery of EIT waves.
\bibitem{Klassen2000} Klassen, A., Aurass, H., Mann, G., \& Thompson, B.~J. 2000, A\&AS,
141, 357 --- EIT-wave speeds $170$--$465\kms$ (mean $271\kms$).
\bibitem{West2011} West, M., Zhukov, A.~N., Dolla, L., \& Rodriguez, L. 2011, ApJ, 730,
122 --- quiet-Sun coronal field $\sim$2--6 G from EIT-wave seismology.
\bibitem{Muhr2011} Muhr, N., Veronig, A.~M., Kienreich, I.~W., Temmer, M., \& Vr\v{s}nak,
B. 2011, ApJ, 739, 89 --- EUV-wave compression ratios 1.05--1.30.
\bibitem{Narukage2002} Narukage, N., et al. 2002, ApJ, 572, L109 --- SXR Moreton-wave
front, fast-mode Mach 1.1--1.3.
\bibitem{Hudson2003} Hudson, H.~S., Khan, J.~I., Lemen, J.~R., Nitta, N.~V., \& Uchida,
Y. 2003, Sol.~Phys., 212, 121.
\bibitem{Cabezas2019} Cabezas, D.~P., et al. 2019, ApJ, 883 (arXiv:1908.03534) ---
chromospheric down-up swing 1.6--4.1$\kms$ (2014 Mar 29 event).
\bibitem{Bala2010} Balasubramaniam, K.~S., Cliver, E.~W., Pevtsov, A., et al. 2010,
ApJ, 723, 587 --- origin of the 2006 December 6 Moreton wave: kinematics
($\sim$850$\kms$ mean, $V_0\simeq1125\kms$, $a=-0.87\;\mathrm{km\,s^{-2}}$, traced
$\sim$120--550 Mm), flare-pulse vs.\ CME-driver analysis.
\bibitem{Balasubramaniam2010} Balasubramaniam, K.~S., Cliver, E.~W., et al. 2010, ApJ,
708, 1639 --- 2.6$\kms$ Doppler swing (2003 Oct 28).
\bibitem{Ballai2005} Ballai, I., Erd\'elyi, R., \& Pint\'er, B. 2005, ApJ, 633, L145 ---
wave energy $\sim$$10^{19}$--$10^{21}$ J from loop-oscillation seismology.
\bibitem{Emslie2012} Emslie, A.~G., et al. 2012, ApJ, 759, 71 --- flare/CME energy
budgets ($10^{29}$--$10^{32}$ erg).
\bibitem{WillsDavey2007} Wills-Davey, M.~J., DeForest, C.~E., \& Stenflo, J.~O. 2007,
ApJ, 664, 556 --- caveat: EIT-wave MHD models requiring unusually low quiet-Sun field.
\bibitem{Aschwanden2005} Aschwanden, M.~J. 2005, \emph{Physics of the Solar Corona}
(Springer) --- coronal scale height, flare-loop $T/n$ limits.
\bibitem{PopescuBraileanu2019} Popescu Br\u{a}ileanu, B., Lukin, V.~S., Khomenko, E., \&
de Vicente, \'A. 2019, A\&A, 627, A25 --- stratified chromospheric waves; exact linear
solution used to validate the SurfaceCode solver.
\end{thebibliography}
\end{document}
